% Anpassung an Landessprache ---------------------------------------------------
\usepackage[ngerman]{babel}

% Umlaute ----------------------------------------------------------------------
%   Umlaute/Sonderzeichen wie äüöß direkt im Quelltext verwenden (CodePage).
% ------------------------------------------------------------------------------
\usepackage[T1]{fontenc} % Erlaubt automatische Trennung von Worten mit Umlauten.
\usepackage{textcomp} % Euro-Zeichen etc.

% Schrift ----------------------------------------------------------------------
\usepackage{lmodern} % bessere Fonts
\usepackage{relsize} % Schriftgröße relativ festlegen

% Tabellen ---------------------------------------------------------------------
\PassOptionsToPackage{table}{xcolor}
\usepackage{tabularx}
\usepackage{multirow}
\usepackage{bigstrut}
% für lange Tabellen
\usepackage{longtable}
\usepackage{array}
\usepackage{ragged2e}
\usepackage{lscape}
\newcolumntype{w}[1]{>{\raggedleft\hspace{0pt}}p{#1}} % Spaltendefinition rechtsbündig mit definierter Breite

% Grafiken ---------------------------------------------------------------------
\usepackage[dvips,final]{graphicx} % Einbinden von JPG-Grafiken ermöglichen
\usepackage{graphics} % keepaspectratio
\usepackage{floatflt} % zum Umfließen von Bildern
\graphicspath{{Bilder/}} % hier liegen die Bilder des Dokuments

\usepackage{tikz} % visualisation
\usetikzlibrary{trees} % directory trees

\usepackage{wrapfig}

% Sonstiges --------------------------------------------------------------------
\usepackage{amssymb} % fuer /checkmark
\usepackage{amsmath,amsfonts} % Befehle aus AMSTeX für mathematische Symbole
\usepackage{enumitem} % anpassbare Enumerates/Itemizes
\usepackage{xspace} % sorgt dafür, dass Leerzeichen hinter parameterlosen Makros nicht als Makroendezeichen interpretiert werden
\usepackage{xparse} % fuer \NewDocumentCommand, ermoeglicht mehrere optionale args
\usepackage{makeidx} % für Index-Ausgabe mit \printindex
\usepackage{acronym} % es werden alle Definitionen aufgelistet | nur benutze : \usepackage[printonlyused]{acronym}

% Einfache Definition der Zeilenabstände und Seitenränder etc.
\usepackage{setspace}
\usepackage{geometry}
% TODO may include
%\usepackage[compact]{titlesec} % Zeilenabstände

% Symbolverzeichnis
\usepackage[intoc]{nomencl}
\let\abbrev\nomenclature
\renewcommand{\nomname}{Abkürzungsverzeichnis}
\setlength{\nomlabelwidth}{.25\hsize}
\renewcommand{\nomlabel}[1]{#1 \dotfill}
\setlength{\nomitemsep}{-\parsep}

\usepackage{url} % URL verlinken, lange URLs umbrechen etc.
\usepackage{chngcntr} % fortlaufendes Durchnummerieren der Fußnoten
\usepackage{ifthen} % bei der Definition eigener Befehle benötigt
\usepackage{todonotes} % definiert u.a. die Befehle \todo und \listoftodos

\usepackage[
    backend=biber, %is default
    style=ieee, % numeric, alphabetic, authoryear, authortitle, science, ieee
    sorting=nyt
]{biblatex} % literaturverzeichnis
\addbibresource{template.bib}

% nty   Sort by name, title, year.
% nyt   Sort by name, year, title.
% nyvt  Sort by name, year, volume, title.
% anyt  Sort by alphabetic label, name, year, title.
% anyvt Sort by alphabetic label, name, year, volume, title.
% ynt   Sort by year, name, title.
% ydnt  Sort by year (descending), name, title.
% none  Do not sort at all. All entries are processed in citation order.


% PDF-Optionen -----------------------------------------------------------------
\usepackage{pdfpages}
\pdfminorversion=5 % erlaubt das Einfügen von pdf-Dateien bis Version 1.7, ohne eine Fehlermeldung zu werfen (keine Garantie für fehlerfreies Einbetten!)
\usepackage[
    bookmarks,
    bookmarksnumbered,
    bookmarksopen=true,
    bookmarksopenlevel=1,
    colorlinks=true,
% diese Farbdefinitionen zeichnen Links im PDF farblich aus
    linkcolor=MyBlau, % einfache interne Verknüpfungen
    anchorcolor=MyBlau,% Ankertext
    citecolor=MyBlau, % Verweise auf Literaturverzeichniseinträge im Text
    filecolor=MyBlau, % Verknüpfungen, die lokale Dateien öffnen
    menucolor=MyBlau, % Acrobat-Menüpunkte
    urlcolor=MyBlau,
% diese Farbdefinitionen sollten für den Druck verwendet werden (alles schwarz)
%linkcolor=black, % einfache interne Verknüpfungen
%anchorcolor=black, % Ankertext
%citecolor=black, % Verweise auf Literaturverzeichniseinträge im Text
%filecolor=black, % Verknüpfungen, die lokale Dateien öffnen
%menucolor=black, % Acrobat-Menüpunkte
%urlcolor=black,
%
%backref, % Quellen werden zurück auf ihre Zitate verlinkt
    pdftex,
    plainpages=false, % zur korrekten Erstellung der Bookmarks
    pdfpagelabels=true, % zur korrekten Erstellung der Bookmarks
    hypertexnames=false, % zur korrekten Erstellung der Bookmarks
    linktoc=all%linktocpage % Seitenzahlen anstatt Text im Inhaltsverzeichnis verlinken
]{hyperref}
% Befehle, die Umlaute ausgeben, führen zu Fehlern, wenn sie hyperref als Optionen übergeben werden
\hypersetup{
    pdftitle={\titel},
    pdfauthor={\autor},
    pdfcreator={\autor},
    pdfsubject={\titel},
    pdfkeywords={\titel},
}


% zum Einbinden von Programmcode -----------------------------------------------
\usepackage{listings}
\usepackage{xcolor}

\definecolor{DEFhellgelb}{rgb}{1,1,0.9}
\definecolor{DEFKeys}{rgb}{0,0,1}
\definecolor{DEFIdentifier}{rgb}{0,0,0}
\definecolor{DEFComments}{rgb}{0,0.5,0}
\definecolor{DEFString}{rgb}{1,0,0}

\definecolor{JAVAred}{rgb}{0.6,0,0}
\definecolor{JAVAgreen}{rgb}{0.25,0.5,0.35}
\definecolor{JAVApurple}{rgb}{0.5,0,0.35}
\definecolor{JAVAdocblue}{rgb}{0.25,0.35,0.75}

\definecolor{forestgreen}{RGB}{34,139,34}
\definecolor{orangered}{RGB}{239,134,64}
\definecolor{darkblue}{rgb}{0.0,0.0,0.6}
\definecolor{gray}{rgb}{0.4,0.4,0.4}

\definecolor{HTMLdarkgray}{rgb}{0.4, 0.4, 0.4}
\definecolor{HTMLocherCode}{rgb}{1, 0.5, 0} % #FF7F00 -> rgb(239, 169, 0)
\definecolor{HTMLblueCode}{rgb}{0, 0, 0.93} % #0000EE -> rgb(0, 0, 238)
\definecolor{HTMLgreenCode}{rgb}{0, 0.6, 0} % #009900 -> rgb(0, 153, 0)

\lstset{
    float=hbp,
    basicstyle=\footnotesize,
    identifierstyle=\color{DEFIdentifier},
    keywordstyle=\color{DEFKeys},
    stringstyle=\color{DEFString},
    commentstyle=\color{DEFComments},
    backgroundcolor=\color{white},
    columns=flexible,
    tabsize=2,
    frame=single,
    extendedchars=true,
    showspaces=false,
    showstringspaces=false,
    numbers=left,
    numberstyle=\tiny,
    breaklines=true,
    breakautoindent=true,
    captionpos=b,
}

\lstdefinestyle{Java}{
    language=Java,
    tabsize=4,
    tabsize=4,
    breakatwhitespace=false,
    resetmargins=true
%
    basicstyle=\footnotesize\ttfamily,
    keywordstyle=\color{JAVApurple}\bfseries,
    stringstyle=\color{JAVAred},
    commentstyle=\color{JAVAgreen},
%
    sensitive=true,
    morecomment=[l]{//},
    morecomment=[s][\color{JAVAdocblue}]{/**}{*/},
    morestring=[b]",
    morestring=[b]',
    morekeywords={
        public,void,package,import,throws,
        abstract,and,array,as,break,case,catch,cfunction,class,clone,const,
        continue,declare,default,do,else,elseif,enddeclare,endfor,endforeach,
        endif,endswitch,endwhile,extends,final,for,foreach,function,global,
        goto,if,implements,interface,instanceof,namespace,new,old_function,or,
        private,protected,public,static,switch,throw,try,use,var,while,xor
        die,echo,empty,exit,eval,include,include_once,isset,list,require,
        require_once,return,print,unset},
}

\lstdefinestyle{XML} {
    basicstyle=\footnotesize,
    columns=flexible,
    tabsize=2,
    frame=single,
    extendedchars=true,
    showspaces=false,
    showstringspaces=false,
    numbers=left,
    numberstyle=\tiny,
    breaklines=true,
    breakautoindent=true,
    captionpos=b,
%
    language=XML,
    extendedchars=true,
    breaklines=true,
    breakatwhitespace=true,
    emph={},
    emphstyle=\color{red},
    basicstyle=\ttfamily,
    columns=fullflexible,
    commentstyle=\color{gray}\upshape,
    morestring=[b]",
    morecomment=[s]{<?}{?>},
    morecomment=[s][\color{forestgreen}]{<!--}{-->},
    keywordstyle=\color{orangered},
    stringstyle=\ttfamily\color{black}\normalfont,
    tagstyle=\color{darkblue}\textbf,
    morekeywords={attribute,xmlns,version,type,release},
}

\lstdefinelanguage{CSS} {
    morekeywords={
        color,background,margin,padding,font,weight,display,position,
        top,left,right,bottom,list,style,border,size,white,space,min,width},
    sensitive=false,
    morecomment=[l]{//},
    morecomment=[s]{/*}{*/},
    morestring=[b]',
    morestring=[b]"
}

\lstdefinelanguage{JavaScript} {
    showstringspaces=false,
    keywordstyle=\ttfamily\bfseries\color{darkblue},
    identifierstyle=\ttfamily,
    stringstyle=\ttfamily\color{forestgreen},
    commentstyle=\color{gray},
    rulecolor=\color{gray},
    xleftmargin=5pt,
    xrightmargin=5pt,
    aboveskip=\bigskipamount,
    belowskip=\bigskipamount,
%
    morekeywords={
        break,const,continue,delete,do,while,export,for,in,function,
        if,else,import,in,instanceOf,label,let,new,return,switch,this,
        throw,try,catch,typeof,var,void,with,yield
    },
    sensitive=false,
    morecomment=[l]{//},
    morecomment=[s]{/*}{*/},
    morestring=[b]",
    morestring=[d]'
}